\documentclass[11pt,a4paper]{article}

\usepackage[utf8]{inputenc}
\usepackage[T1]{fontenc}
\usepackage[french]{babel}
\usepackage[a4paper,margin=1.8cm]{geometry}
\usepackage{amsmath,amssymb,amsthm,mathtools}
\usepackage{lmodern}
\usepackage{xcolor}
\usepackage{enumitem}
\usepackage{array}
\usepackage{booktabs}
\usepackage{longtable}
\usepackage{hyperref}
\usepackage{fancyhdr}
\usepackage{titlesec}
\usepackage{tcolorbox}
\usepackage{multicol}
\usepackage{tikz}
\usepackage{physics}

\hypersetup{
    colorlinks=true,
    linkcolor=black,
    urlcolor=blue
}

\setlength{\parindent}{0pt}
\setlength{\parskip}{0.45em}

\pagestyle{fancy}
\fancyhf{}
\fancyhead[L]{Fiche de synthèse -- Physique-Chimie}
\fancyhead[R]{Première \& Terminale générale}
\fancyfoot[C]{\thepage}

\titleformat{\section}{\Large\bfseries}{\thesection.}{0.5em}{}
\titleformat{\subsection}{\large\bfseries}{\thesubsection.}{0.5em}{}
\titleformat{\subsubsection}{\normalsize\bfseries}{\thesubsubsection.}{0.5em}{}

\newcommand{\R}{\mathbb{R}}
\newcommand{\dd}{\,\mathrm{d}}
\newcommand{\e}{\mathrm{e}}
\newcommand{\jj}{\mathrm{j}}
\newcommand{\ii}{\mathrm{i}}

\newtcolorbox{idee}{
colback=blue!4,
colframe=blue!45!black,
title=Idée pédagogique
}

\newtcolorbox{erreurs}{
colback=red!3,
colframe=red!50!black,
title=Erreurs classiques des élèves
}

\newtcolorbox{demo}{
colback=green!3,
colframe=green!45!black,
title=Démonstration importante
}

\newtcolorbox{modele}{
colback=yellow!7,
colframe=orange!60!black,
title=Comprendre le modèle
}

\newtcolorbox{bilan}{
colback=gray!8,
colframe=black!60,
title=À retenir absolument
}

\newtcolorbox{sens}{
colback=cyan!4,
colframe=cyan!50!black,
title=Sens concret de la notion
}

\newtcolorbox{imagem}{
colback=violet!4,
colframe=violet!50!black,
title=Image mentale utile
}

\newtcolorbox{fausseintuition}{
colback=pink!4,
colframe=pink!60!black,
title=Fausse intuition fréquente
}

\begin{document}

\begin{center}
    {\LARGE \textbf{Fiche de synthèse pour un professeur de physique-chimie}}\\[0.4em]
    {\large \textbf{Première \& Terminale générale -- spécialité}}\\[0.4em]
    Notions clés, sens physique, démonstrations importantes, modèles, difficultés conceptuelles et erreurs classiques.
\end{center}

\vspace{0.8em}

\tableofcontents

\newpage

\section{Objectif général du document}

Ce document est conçu comme une fiche de synthèse pour un professeur de physique-chimie au lycée.
Il ne vise pas seulement à rappeler les notions du programme, mais à mettre en évidence :
\begin{itemize}[label=--]
    \item les \textbf{concepts fondamentaux} à faire construire ;
    \item le \textbf{sens concret} des grandeurs manipulées ;
    \item les \textbf{modèles} utilisés et leur domaine de validité ;
    \item les \textbf{démonstrations importantes} qui structurent la compréhension ;
    \item les \textbf{erreurs classiques des élèves} ;
    \item les \textbf{images mentales utiles} pour enseigner sans trahir les concepts.
\end{itemize}

\begin{bilan}
Une difficulté majeure en physique-chimie au lycée est que les élèves apprennent souvent des formules
sans comprendre ce qu'elles signifient physiquement.
Le rôle du professeur est donc de faire constamment le lien entre :
\begin{itemize}
    \item l'observation ;
    \item la grandeur mesurée ;
    \item le modèle choisi ;
    \item l'écriture mathématique ;
    \item l'interprétation physique.
\end{itemize}
\end{bilan}

\begin{idee}
Pour chaque notion, il est utile de faire expliciter aux élèves :
\begin{enumerate}
    \item ce que la grandeur mesure réellement ;
    \item à quoi elle correspond concrètement ;
    \item comment on la distingue d'autres notions voisines ;
    \item dans quel cadre le modèle est valable.
\end{enumerate}
\end{idee}

\section{Le rôle central des modèles}

Un modèle physique n'est pas la réalité.
C'est une représentation simplifiée construite pour décrire, prévoir ou interpréter une situation.

\begin{modele}
Un bon modèle n'est pas un modèle \og vrai \fg\ au sens absolu.
C'est un modèle :
\begin{itemize}
    \item assez simple pour être exploitable ;
    \item assez fidèle pour rendre compte du phénomène étudié ;
    \item adapté à l'échelle considérée.
\end{itemize}
\end{modele}

\subsection{Pourquoi utilise-t-on le modèle de Newton au lycée ?}

\begin{modele}
Le modèle newtonien est utilisé parce qu'il relie directement :
\[
\sum \vec F = m\vec a
\]
et permet donc :
\begin{itemize}
    \item d'expliquer qualitativement un mouvement ;
    \item de prévoir quantitativement son évolution ;
    \item de faire le lien entre une situation concrète et une écriture mathématique simple.
\end{itemize}

Il est adapté aux situations usuelles du lycée :
\begin{itemize}
    \item vitesses faibles devant celle de la lumière ;
    \item masses et distances ordinaires ;
    \item référentiels terrestres assimilés à des référentiels galiléens dans de nombreux cas.
\end{itemize}
\end{modele}

\begin{idee}
Faire comprendre aux élèves qu'on ne prétend pas décrire toute la réalité dans ses détails,
mais seulement choisir un cadre suffisamment bon pour répondre à la question posée.
\end{idee}

\section{Donner du sens concret aux notions physiques}

Une notion physique n'est bien comprise que si l'élève peut répondre à trois questions :
\begin{enumerate}
    \item Qu'est-ce que cette grandeur mesure réellement ?
    \item À quoi correspond-elle concrètement ?
    \item Quelle erreur d'interprétation faut-il éviter ?
\end{enumerate}

\begin{idee}
Une image mentale est très utile pour enseigner, mais elle ne remplace jamais la définition rigoureuse.
Il faut souvent dire explicitement :
\textit{``Ceci est une image pour comprendre, pas la définition exacte.''}
\end{idee}

\newpage

\section{Première générale -- Mouvement et interactions}

\subsection{Position}

\begin{sens}
La position indique où se trouve un système par rapport à un repère choisi.
Ce n'est jamais une propriété absolue : elle dépend du référentiel et de l'origine choisie.
\end{sens}

\begin{imagem}
La position, c'est l'\textbf{adresse} de l'objet dans l'espace.
\end{imagem}

\begin{fausseintuition}
Les élèves croient souvent qu'un objet a une position \og absolue \fg.
Il faut rappeler qu'on ne peut parler de position que relativement à un repère.
\end{fausseintuition}

\subsection{Référentiel}

\begin{sens}
Le référentiel est le cadre dans lequel on décrit le mouvement.
Il fixe à la fois un repère d'espace et une mesure du temps.
\end{sens}

\begin{imagem}
Le référentiel, c'est le \textbf{point de vue choisi pour raconter le mouvement}.
\end{imagem}

\begin{erreurs}
\begin{itemize}
    \item oublier de préciser le référentiel ;
    \item croire qu'un mouvement existe indépendamment de tout observateur ;
    \item mélanger plusieurs référentiels sans le dire.
\end{itemize}
\end{erreurs}

\subsection{Vitesse}

\begin{sens}
La vitesse mesure à quel point la position change rapidement.
Sous forme vectorielle, elle donne aussi la direction du mouvement.
\end{sens}

\begin{imagem}
La vitesse répond à la question :
\textit{``à quelle vitesse et dans quelle direction l'objet change-t-il de place ?''}
\end{imagem}

\begin{fausseintuition}
Les élèves réduisent souvent la vitesse à un simple nombre positif.
Il faut insister sur le fait que le vecteur vitesse a une direction et un sens.
\end{fausseintuition}

\subsection{Accélération}

\begin{sens}
L'accélération mesure la vitesse de variation de la vitesse.
Elle apparaît dès que la valeur de la vitesse change ou dès que sa direction change.
\end{sens}

\begin{imagem}
L'accélération est le \textbf{thermomètre du changement de mouvement}.
\end{imagem}

\begin{fausseintuition}
Les élèves pensent souvent qu'un objet très rapide a forcément une grande accélération.
Un mouvement très rapide peut être uniforme et avoir une accélération nulle.
\end{fausseintuition}

\subsection{Forces}

\begin{sens}
Une force modélise une action mécanique exercée sur un système.
Elle permet d'expliquer une déformation ou une modification du mouvement.
\end{sens}

\begin{imagem}
Une force est une \textbf{influence extérieure} qui pousse, tire, freine, dévie ou comprime.
\end{imagem}

\begin{fausseintuition}
Beaucoup d'élèves pensent qu'il faut une force pour maintenir un mouvement.
Dans le modèle de Newton, une force est nécessaire pour changer le mouvement, pas pour conserver une vitesse constante.
\end{fausseintuition}

\subsection{Principe d'inertie}

\begin{bilan}
Dans un référentiel galiléen, si la résultante des forces extérieures est nulle,
le vecteur vitesse reste constant.
\end{bilan}

\begin{demo}
Si
\[
\sum \vec F = \vec 0,
\]
alors, dans le cadre du modèle newtonien,
\[
m\vec a=\vec 0.
\]
Comme \(m>0\), on en déduit
\[
\vec a=\vec 0.
\]
Donc
\[
\dv{\vec v}{t}=\vec 0,
\]
ce qui implique que \(\vec v\) est constant.
Le mouvement est donc rectiligne uniforme, ou le système reste au repos.
\end{demo}

\begin{modele}
Le principe d'inertie est une loi de non-changement du mouvement.
Il permet de rompre avec l'idée intuitive selon laquelle il faudrait une force pour qu'un objet continue à avancer.
\end{modele}

\begin{erreurs}
\begin{itemize}
    \item croire que l'absence de force implique l'absence de mouvement ;
    \item croire qu'un mouvement uniforme nécessite une force motrice ;
    \item confondre vitesse nulle et accélération nulle.
\end{itemize}
\end{erreurs}

\subsection{Deuxième loi de Newton}

\begin{bilan}
\[
\sum \vec F = m\vec a
\]
pour un système de masse constante dans un référentiel galiléen.
\end{bilan}

\begin{sens}
Cette loi relie les causes mécaniques modélisées par les forces
et l'effet observé, à savoir l'accélération.
\end{sens}

\begin{imagem}
Les forces jouent le rôle des \textbf{commandes du changement de mouvement}.
Elles ne pilotent pas directement la vitesse, mais sa variation.
\end{imagem}

\begin{demo}
Dans un repère orthonormé \((O,\vec i,\vec j)\), on écrit :
\[
\sum \vec F = \left(\sum F_x\right)\vec i + \left(\sum F_y\right)\vec j
\]
et
\[
m\vec a = ma_x\vec i + ma_y\vec j.
\]
Comme deux vecteurs sont égaux si et seulement si leurs composantes sont égales dans une même base,
on obtient :
\[
\sum F_x = ma_x
\qquad \text{et} \qquad
\sum F_y = ma_y.
\]
Cette étape justifie les équations projetées.
\end{demo}

\begin{erreurs}
\begin{itemize}
    \item oublier que la loi est vectorielle ;
    \item sommer algébriquement sans projeter ;
    \item confondre \(\vec v\) et \(\vec a\) ;
    \item mal gérer les signes selon l'axe choisi.
\end{itemize}
\end{erreurs}

\subsection{Poids}

\begin{sens}
Le poids modélise l'action gravitationnelle exercée par la Terre sur un système proche de sa surface.
\end{sens}

\begin{imagem}
Le poids, c'est l'\textbf{appel de la Terre} vers son centre.
\end{imagem}

\begin{bilan}
\[
\vec P = m\vec g.
\]
\end{bilan}

\begin{fausseintuition}
Les élèves confondent souvent masse et poids.
La masse mesure l'inertie, le poids est une force.
\end{fausseintuition}

\subsection{Chute libre et tir parabolique}

\begin{modele}
En chute libre au lycée, on suppose que seule la pesanteur agit et que les frottements de l'air sont négligés.
Il faut toujours expliciter cette hypothèse.
\end{modele}

\begin{demo}
Si seule la pesanteur agit :
\[
m\vec a = m\vec g,
\]
donc
\[
\vec a = \vec g.
\]
Dans un repère où \(x\) est horizontal et \(z\) vertical vers le haut :
\[
a_x=0,\qquad a_z=-g.
\]
En intégrant :
\[
v_x(t)=v_{0x},\qquad v_z(t)=v_{0z}-gt
\]
puis
\[
x(t)=x_0+v_{0x}t,\qquad z(t)=z_0+v_{0z}t-\frac12gt^2.
\]
Le mouvement est uniforme horizontalement et uniformément varié verticalement.
\end{demo}

\begin{erreurs}
\begin{itemize}
    \item croire que la vitesse horizontale diminue forcément ;
    \item se tromper sur le signe de \(g\) ;
    \item oublier qu'une trajectoire parabolique résulte de la composition de deux mouvements.
\end{itemize}
\end{erreurs}

\newpage

\section{Première générale -- Énergie : conversions et transferts}

\subsection{Travail d'une force}

\begin{sens}
Le travail d'une force mesure le transfert d'énergie associé à cette force
lorsqu'un système se déplace.
\end{sens}

\begin{imagem}
Le travail est la \textbf{part d'action utile} d'une force le long du déplacement.
\end{imagem}

\begin{bilan}
Pour une force constante :
\[
W_{A\to B}(\vec F)=\vec F\cdot \overrightarrow{AB}=F\,AB\cos\theta.
\]
\end{bilan}

\begin{demo}
Par définition, le travail d'une force constante sur un déplacement \(\overrightarrow{AB}\) vaut :
\[
W_{A\to B}(\vec F)=\vec F\cdot \overrightarrow{AB}.
\]
Or, par définition du produit scalaire :
\[
\vec F\cdot \overrightarrow{AB}=F\,AB\cos\theta.
\]
On obtient immédiatement :
\begin{itemize}
    \item un travail positif si la force aide le déplacement ;
    \item nul si elle lui est perpendiculaire ;
    \item négatif si elle s'y oppose.
\end{itemize}
\end{demo}

\begin{erreurs}
\begin{itemize}
    \item confondre force et travail ;
    \item oublier le rôle de l'angle ;
    \item mal interpréter un travail négatif.
\end{itemize}
\end{erreurs}

\subsection{Énergie}

\begin{sens}
L'énergie est une grandeur qui mesure la capacité d'un système à produire une transformation :
mettre en mouvement, chauffer, déformer, alimenter un appareil, rayonner, etc.
\end{sens}

\begin{imagem}
L'énergie peut être vue comme une \textbf{monnaie universelle des transformations physiques}.
\end{imagem}

\begin{fausseintuition}
Les élèves imaginent souvent l'énergie comme une substance contenue dans l'objet.
Il est plus juste de dire qu'il s'agit d'une grandeur de bilan permettant de décrire et quantifier des transformations.
\end{fausseintuition}

\subsection{Énergie cinétique}

\begin{sens}
L'énergie cinétique est l'énergie associée au mouvement.
Plus un système est massif et rapide, plus elle est grande.
\end{sens}

\begin{imagem}
L'énergie cinétique représente ce que le mouvement peut produire comme effet :
choc, freinage, déformation, échauffement.
\end{imagem}

\begin{bilan}
\[
E_c=\frac12 mv^2.
\]
\end{bilan}

\begin{fausseintuition}
Les élèves confondent souvent énergie cinétique et quantité de mouvement.
L'énergie cinétique est scalaire, la quantité de mouvement est vectorielle.
\end{fausseintuition}

\subsection{Théorème de l'énergie cinétique}

\begin{bilan}
\[
\Delta E_c = \sum W(\vec F).
\]
\end{bilan}

\begin{demo}
Dans un cas unidimensionnel, on part de
\[
\sum F = ma.
\]
Or
\[
a=\dv{v}{t}
\qquad \text{et} \qquad
\dd x = v\dd t.
\]
Donc
\[
\sum F\,\dd x = m\dv{v}{t}v\dd t = mv\dd v.
\]
Ainsi :
\[
\sum F\,\dd x = \dd\left(\frac12 mv^2\right).
\]
Après intégration entre \(A\) et \(B\), on obtient :
\[
\sum W_{A\to B}(\vec F)=E_{c,B}-E_{c,A}.
\]
\end{demo}

\begin{modele}
Ce théorème montre la cohérence entre dynamique et énergétique.
Il ne s'agit pas d'une formule indépendante, mais d'une autre manière de lire la même physique.
\end{modele}

\subsection{Énergie potentielle de pesanteur}

\begin{sens}
L'énergie potentielle de pesanteur est l'énergie associée à la position d'un système dans le champ de pesanteur.
\end{sens}

\begin{imagem}
C'est une \textbf{réserve de possibilité de descente} :
un objet placé en hauteur peut convertir cette énergie en mouvement.
\end{imagem}

\begin{bilan}
\[
E_p = mgz
\]
à une constante de référence près.
\end{bilan}

\begin{fausseintuition}
Les élèves pensent souvent que l'énergie potentielle est absolue.
Or elle dépend du niveau de référence choisi.
Ce sont surtout les variations d'énergie potentielle qui ont un sens physique direct.
\end{fausseintuition}

\subsection{Énergie mécanique}

\begin{sens}
L'énergie mécanique est la somme de l'énergie cinétique et de l'énergie potentielle.
Elle représente le budget énergétique lié à la position et au mouvement.
\end{sens}

\begin{imagem}
Lors d'une chute sans frottement, on échange de la \og hauteur \fg\ contre de la \og vitesse \fg.
L'énergie mécanique est alors le budget total qui se redistribue.
\end{imagem}

\begin{bilan}
\[
E_m = E_c + E_p.
\]
\end{bilan}

\begin{demo}
Si seul le poids agit :
\[
W_{A\to B}(\vec P)=mg(z_A-z_B).
\]
Par le théorème de l'énergie cinétique :
\[
E_{c,B}-E_{c,A}=mg(z_A-z_B).
\]
Donc
\[
E_{c,B}+mgz_B=E_{c,A}+mgz_A,
\]
ce qui montre que
\[
E_m=E_c+E_p
\]
reste constante.
\end{demo}

\begin{erreurs}
\begin{itemize}
    \item croire que l'énergie mécanique se conserve toujours ;
    \item oublier le rôle des frottements ;
    \item confondre conservation de l'énergie mécanique et conservation de l'énergie totale.
\end{itemize}
\end{erreurs}

\subsection{Puissance}

\begin{sens}
La puissance mesure la vitesse à laquelle de l'énergie est transférée ou convertie.
\end{sens}

\begin{imagem}
La puissance, c'est le \textbf{rythme énergétique}.
Deux systèmes peuvent consommer la même énergie totale, mais pas avec la même rapidité.
\end{imagem}

\begin{bilan}
\[
P=\frac{\dd W}{\dd t}
\qquad \text{et} \qquad
P=\vec F\cdot \vec v.
\]
\end{bilan}

\begin{fausseintuition}
Les élèves confondent souvent énergie et puissance.
L'énergie est une quantité totale, la puissance est une quantité par unité de temps.
\end{fausseintuition}

\newpage

\section{Première générale -- Ondes et signaux}

\subsection{Onde}

\begin{sens}
Une onde est la propagation d'une perturbation de proche en proche.
Elle transporte de l'énergie et de l'information, sans transport global de matière.
\end{sens}

\begin{imagem}
Une onde, c'est la \textbf{forme d'une perturbation qui se déplace},
pas un morceau de matière qui voyage.
\end{imagem}

\begin{fausseintuition}
Les élèves pensent souvent que la matière est emportée avec l'onde.
Dans une corde, ce n'est pas la corde entière qui part avec la perturbation.
\end{fausseintuition}

\subsection{Célérité}

\begin{sens}
La célérité mesure la vitesse de propagation de l'onde dans un milieu donné.
\end{sens}

\begin{imagem}
La célérité est la \textbf{vitesse de voyage de la perturbation}.
\end{imagem}

\begin{fausseintuition}
Il ne faut pas la confondre avec la vitesse locale des points du milieu.
\end{fausseintuition}

\subsection{Période et fréquence}

\begin{sens}
La période \(T\) est la durée d'un motif temporel élémentaire.
La fréquence \(f\) mesure le nombre de répétitions par seconde.
\end{sens}

\begin{imagem}
La fréquence est le \textbf{rythme de répétition} d'un phénomène périodique.
\end{imagem}

\begin{bilan}
\[
f=\frac{1}{T}.
\]
\end{bilan}

\begin{demo}
Si un phénomène se reproduit identiquement toutes les \(T\) secondes,
alors en une seconde il effectue
\[
f=\frac{1}{T}
\]
périodes.
\end{demo}

\begin{erreurs}
\begin{itemize}
    \item confondre période et fréquence ;
    \item oublier les unités ;
    \item confondre fréquence et vitesse de propagation.
\end{itemize}
\end{erreurs}

\subsection{Longueur d'onde}

\begin{sens}
La longueur d'onde est la distance parcourue par l'onde pendant une période.
C'est aussi la distance entre deux points successifs en phase.
\end{sens}

\begin{imagem}
C'est la \textbf{taille spatiale d'un motif élémentaire} de l'onde.
\end{imagem}

\begin{fausseintuition}
Les élèves mélangent souvent longueur d'onde et amplitude.
L'une décrit un espacement, l'autre l'intensité de la perturbation.
\end{fausseintuition}

\subsection{Lumière et spectres}

\begin{sens}
L'étude des spectres permet d'identifier des espèces ou des phénomènes à partir de l'interaction entre la matière et le rayonnement.
\end{sens}

\begin{imagem}
Un spectre est une \textbf{empreinte lumineuse} propre à une situation physique ou chimique.
\end{imagem}

\begin{erreurs}
\begin{itemize}
    \item lire un spectre sans regarder les axes ;
    \item identifier une espèce à partir d'un seul indice ;
    \item oublier qu'un spectre est un résultat expérimental à interpréter.
\end{itemize}
\end{erreurs}

\newpage

\section{Première générale -- Constitution et transformations de la matière}

\subsection{Quantité de matière}

\begin{sens}
La quantité de matière sert à compter les entités chimiques à l'échelle macroscopique.
\end{sens}

\begin{imagem}
C'est une manière pratique de compter des particules en quantités immenses,
comme on compterait des feuilles par paquets plutôt qu'une par une.
\end{imagem}

\begin{bilan}
\[
n=\frac{m}{M}.
\]
\end{bilan}

\begin{fausseintuition}
Les élèves confondent souvent quantité de matière et masse.
Deux échantillons de même masse peuvent correspondre à des quantités de matière différentes.
\end{fausseintuition}

\subsection{Concentration}

\begin{sens}
La concentration mesure la quantité d'espèce chimique présente dans un volume donné.
\end{sens}

\begin{imagem}
La concentration correspond à une \textbf{densité de présence} des entités dans la solution.
\end{imagem}

\begin{bilan}
\[
c=\frac{n}{V}.
\]
\end{bilan}

\begin{fausseintuition}
Une solution concentrée n'est pas seulement une solution qui contient \og beaucoup de matière \fg,
mais une solution qui contient beaucoup d'espèce dissoute \textbf{par unité de volume}.
\end{fausseintuition}

\subsection{Dilution}

\begin{modele}
Dans une dilution, on ajoute du solvant sans ajouter de soluté.
La quantité de soluté est donc conservée.
\end{modele}

\begin{demo}
Avant dilution :
\[
n_i=c_iV_i.
\]
Après dilution :
\[
n_f=c_fV_f.
\]
Comme la quantité de soluté ne change pas :
\[
n_i=n_f.
\]
Donc
\[
c_iV_i=c_fV_f.
\]
\end{demo}

\begin{erreurs}
\begin{itemize}
    \item oublier que le volume doit être exprimé dans des unités cohérentes ;
    \item croire qu'on conserve la concentration ;
    \item utiliser la formule sans comprendre la conservation du soluté.
\end{itemize}
\end{erreurs}

\subsection{Transformation chimique et avancement}

\begin{sens}
L'avancement mesure l'état de progression d'une transformation chimique.
\end{sens}

\begin{imagem}
L'avancement est le \textbf{compteur de progression} de la réaction.
\end{imagem}

\begin{bilan}
Pour
\[
aA+bB \longrightarrow cC+dD,
\]
on écrit :
\[
n_A=n_{A,0}-ax,\quad
n_B=n_{B,0}-bx,\quad
n_C=n_{C,0}+cx,\quad
n_D=n_{D,0}+dx.
\]
\end{bilan}

\begin{demo}
Lorsque la réaction avance de \(x\) mole, elle consomme \(ax\) mole de \(A\) et \(bx\) mole de \(B\),
et elle forme \(cx\) mole de \(C\) et \(dx\) mole de \(D\).
On obtient donc directement les expressions des quantités finales :
\[
n_f = n_i + \Delta n.
\]
\end{demo}

\begin{erreurs}
\begin{itemize}
    \item oublier les coefficients stœchiométriques ;
    \item mettre le mauvais signe ;
    \item confondre avancement et quantité de matière d'une espèce.
\end{itemize}
\end{erreurs}

\subsection{Acide et base}

\begin{sens}
Un acide est une espèce capable de céder un proton \(H^+\),
une base est une espèce capable de capter un proton.
\end{sens}

\begin{imagem}
C'est un \textbf{jeu de transmission du proton} :
l'acide donne, la base reçoit.
\end{imagem}

\begin{fausseintuition}
Les élèves réduisent parfois la notion de base à la seule présence de l'ion hydroxyde.
Il faut revenir à la définition générale par transfert de proton.
\end{fausseintuition}

\subsection{Oxydant et réducteur}

\begin{sens}
Un oxydant capte des électrons, un réducteur cède des électrons.
\end{sens}

\begin{imagem}
C'est un \textbf{jeu de transmission d'électrons}.
\end{imagem}

\begin{fausseintuition}
Erreur très fréquente :
\og l'oxydant s'oxyde \fg.
Il faut rappeler :
l'oxydant oxyde l'autre espèce, donc il se réduit.
\end{fausseintuition}

\begin{erreurs}
\begin{itemize}
    \item oublier d'équilibrer les charges ;
    \item ne pas distinguer demi-équations et équation globale ;
    \item confondre charge électrique et nombre d'atomes.
\end{itemize}
\end{erreurs}

\newpage

\section{Terminale générale -- Constitution et transformations de la matière}

\subsection{Cinétique chimique}

\begin{sens}
La cinétique chimique étudie la vitesse d'évolution des transformations.
Elle ne dit pas seulement ce qui est possible, mais à quel rythme cela se produit.
\end{sens}

\begin{imagem}
La cinétique répond à la question :
\textit{``la réaction a lieu, d'accord, mais est-ce maintenant, lentement, très vite ?''}
\end{imagem}

\begin{fausseintuition}
Les élèves confondent souvent réaction totale et réaction rapide.
Une réaction peut être quasi totale tout en étant lente.
\end{fausseintuition}

\begin{modele}
Ce chapitre est très important pour distinguer :
\begin{itemize}
    \item l'état final ;
    \item le temps nécessaire pour l'atteindre.
\end{itemize}
\end{modele}

\subsection{Loi exponentielle}

\begin{bilan}
Si une grandeur \(N(t)\) vérifie
\[
\dv{N}{t}=-\lambda N,
\]
alors
\[
N(t)=N_0\e^{-\lambda t}.
\]
\end{bilan}

\begin{demo}
On écrit :
\[
\frac{1}{N}\dv{N}{t}=-\lambda.
\]
Donc
\[
\frac{\dd N}{N}=-\lambda \dd t.
\]
En intégrant :
\[
\ln N = -\lambda t + C.
\]
D'où
\[
N(t)=A\e^{-\lambda t}.
\]
Avec \(N(0)=N_0\), on obtient
\[
N(t)=N_0\e^{-\lambda t}.
\]
\end{demo}

\begin{idee}
Cette démonstration est très formatrice :
elle relie modélisation, équation différentielle, mathématiques et lecture physique du paramètre \(\lambda\).
\end{idee}

\subsection{Titrage}

\begin{sens}
Le titrage permet de déterminer quantitativement une concentration inconnue
à partir d'une réaction de stœchiométrie connue.
\end{sens}

\begin{imagem}
Le titrage revient à \textbf{compter indirectement} une quantité inconnue
en la faisant réagir avec une quantité maîtrisée.
\end{imagem}

\begin{bilan}
À l'équivalence, les réactifs ont été introduits dans les proportions stœchiométriques :
\[
\frac{n_A}{a}=\frac{n_B}{b}.
\]
\end{bilan}

\begin{erreurs}
\begin{itemize}
    \item utiliser la relation sans expliciter la réaction support ;
    \item confondre espèce titrée et espèce titrante ;
    \item mal lire le point d'équivalence sur une courbe.
\end{itemize}
\end{erreurs}

\subsection{Pile et électrolyse}

\begin{sens}
Une pile convertit de l'énergie chimique en énergie électrique.
Une électrolyse réalise l'inverse : elle force une transformation chimique grâce à un apport électrique.
\end{sens}

\begin{imagem}
On peut voir la pile comme un \textbf{moteur chimique du courant},
et l'électrolyse comme une \textbf{machine électrique qui pousse la chimie}.
\end{imagem}

\begin{bilan}
\[
Q=I\Delta t
\qquad \text{et} \qquad
Q=n(e^-)F.
\]
\end{bilan}

\begin{fausseintuition}
Les élèves confondent souvent sens du courant et sens du déplacement des électrons.
Il faut distinguer rigoureusement les deux.
\end{fausseintuition}

\newpage

\section{Terminale générale -- Mouvement et interactions}

\subsection{Champ électrique}

\begin{sens}
Le champ électrique décrit comment une charge électrique placée en un point serait sollicitée électriquement.
\end{sens}

\begin{imagem}
Le champ électrique peut être vu comme une \textbf{carte d'influence électrique de l'espace}.
\end{imagem}

\begin{bilan}
\[
\vec F = q\vec E.
\]
\end{bilan}

\begin{fausseintuition}
Les élèves pensent parfois que la force a toujours le même sens que le champ.
C'est vrai pour une charge positive, faux pour une charge négative.
\end{fausseintuition}

\subsection{Mouvement circulaire}

\begin{sens}
Un mouvement circulaire est un mouvement dont la trajectoire est un cercle.
Même si la norme de la vitesse reste constante, le vecteur vitesse change de direction.
\end{sens}

\begin{imagem}
Un objet en mouvement circulaire \textbf{tourne sans cesse sa vitesse}.
Il faut donc une accélération qui le ramène vers le centre.
\end{imagem}

\begin{fausseintuition}
Les élèves croient souvent :
\og vitesse constante donc accélération nulle \fg.
C'est faux si la direction change.
\end{fausseintuition}

\subsection{Satellites et gravitation}

\begin{sens}
Le mouvement orbital résulte de l'interaction gravitationnelle entre deux corps.
La gravitation fournit ici la force centripète nécessaire.
\end{sens}

\begin{imagem}
Un satellite est en \textbf{chute permanente autour de la Terre}.
Il tombe sans cesse, mais la courbure de sa trajectoire lui fait manquer le sol.
\end{imagem}

\begin{demo}
Pour un satellite de masse \(m\) en orbite circulaire de rayon \(r\) autour d'un astre de masse \(M\),
la gravitation fournit la force centripète :
\[
\frac{GMm}{r^2}=m\frac{v^2}{r}.
\]
En simplifiant par \(m\), on obtient :
\[
\frac{GM}{r^2}=\frac{v^2}{r}
\]
puis
\[
v=\sqrt{\frac{GM}{r}}.
\]
\end{demo}

\begin{erreurs}
\begin{itemize}
    \item confondre rayon de l'orbite et rayon de la planète ;
    \item oublier que la gravitation décroît comme \(1/r^2\) ;
    \item croire qu'il existe une force \og centrifuge \fg\ réelle agissant sur le satellite dans le référentiel d'étude ordinaire.
\end{itemize}
\end{erreurs}

\newpage

\section{Terminale générale -- Énergie : conversions et transferts}

\subsection{Température}

\begin{sens}
La température caractérise l'état thermique d'un système.
À l'échelle microscopique, elle est reliée au degré d'agitation des constituants.
\end{sens}

\begin{imagem}
La température renseigne sur \textbf{l'intensité de l'agitation microscopique},
pas directement sur une quantité totale d'énergie.
\end{imagem}

\begin{fausseintuition}
Les élèves confondent souvent température et chaleur.
Un objet petit et très chaud n'a pas forcément plus d'énergie thermique qu'un objet grand et tiède.
\end{fausseintuition}

\subsection{Transfert thermique}

\begin{sens}
Un transfert thermique est un transfert d'énergie dû à une différence de température.
\end{sens}

\begin{imagem}
La chaleur n'est pas un contenu stocké :
c'est \textbf{ce qui passe} d'un corps à un autre quand ils ne sont pas à la même température.
\end{imagem}

\begin{fausseintuition}
Dire qu'un corps \og contient de la chaleur \fg\ est trompeur.
Le mot chaleur doit plutôt désigner un transfert.
\end{fausseintuition}

\subsection{Bilan thermique}

\begin{bilan}
Dans un système isolé :
\[
Q_{\text{total}}=0.
\]
Pour deux corps :
\[
Q_1+Q_2=0.
\]
\end{bilan}

\begin{demo}
Si le système global n'échange pas d'énergie avec l'extérieur,
la somme algébrique des transferts thermiques internes est nulle.
Pour deux corps, cela donne directement
\[
Q_1+Q_2=0.
\]
Avec
\[
Q=mc(T_f-T_i),
\]
on obtient une équation permettant de déterminer la température finale.
\end{demo}

\begin{erreurs}
\begin{itemize}
    \item erreurs de signe ;
    \item confusion entre chaleur reçue et température finale ;
    \item oublier qu'un bilan doit être écrit sur le système total.
\end{itemize}
\end{erreurs}

\subsection{Charge électrique}

\begin{sens}
La charge électrique est une propriété de la matière qui rend compte des interactions électriques.
\end{sens}

\begin{imagem}
La charge est ce qui permet à la matière de \textbf{réagir électriquement} :
attirer, repousser, créer un champ.
\end{imagem}

\begin{fausseintuition}
Les élèves confondent souvent charge et courant.
La charge est une grandeur portée par la matière, le courant est un débit de charge.
\end{fausseintuition}

\subsection{Courant électrique}

\begin{sens}
Le courant électrique mesure le débit de charge traversant une section de circuit.
\end{sens}

\begin{imagem}
Le courant, c'est un \textbf{débit} :
comme dans une rivière, ce qui compte est la quantité qui passe par seconde.
\end{imagem}

\begin{bilan}
\[
I=\frac{\Delta Q}{\Delta t}.
\]
\end{bilan}

\begin{fausseintuition}
Le courant n'est pas une réserve stockée dans le fil.
C'est un flux, un rythme de passage.
\end{fausseintuition}

\subsection{Tension électrique}

\begin{sens}
La tension traduit une différence de potentiel électrique entre deux points.
\end{sens}

\begin{imagem}
La tension est une \textbf{différence de niveau énergétique par unité de charge}.
\end{imagem}

\begin{fausseintuition}
La tension ne circule pas.
Elle compare deux points.
Le courant, lui, correspond au passage de charge.
\end{fausseintuition}

\subsection{Résistance}

\begin{sens}
La résistance traduit l'opposition d'un dipôle au passage du courant.
\end{sens}

\begin{imagem}
Une résistance, c'est comme un \textbf{passage étroit} dans une canalisation.
\end{imagem}

\begin{fausseintuition}
La résistance ne \og détruit \fg\ pas le courant.
Elle limite le courant et participe à la conversion de l'énergie électrique, souvent en énergie thermique.
\end{fausseintuition}

\subsection{Puissance électrique}

\begin{sens}
La puissance électrique mesure la vitesse de transfert ou de conversion de l'énergie électrique.
\end{sens}

\begin{imagem}
La puissance, c'est le \textbf{rythme de consommation ou de fourniture d'énergie}.
\end{imagem}

\begin{fausseintuition}
Les élèves confondent souvent puissance et énergie.
La puissance est une énergie par unité de temps.
\end{fausseintuition}

\subsection{Circuit RC}

\begin{modele}
Le dipôle RC est un modèle simple mais très formateur :
il permet d'étudier un régime transitoire, une loi exponentielle,
et le rôle d'une constante de temps.
\end{modele}

\begin{bilan}
Pour la charge d'un condensateur :
\[
RC\dv{u_C}{t}+u_C=E
\]
et
\[
u_C(t)=E\left(1-\e^{-t/RC}\right).
\]
\end{bilan}

\begin{demo}
Dans un circuit série :
\[
u_R+u_C=E.
\]
Or
\[
u_R=Ri
\qquad \text{et} \qquad
i=C\dv{u_C}{t}.
\]
Donc
\[
RC\dv{u_C}{t}+u_C=E.
\]
La solution de cette équation différentielle est :
\[
u_C(t)=E\left(1-\e^{-t/RC}\right).
\]
\end{demo}

\begin{imagem}
La constante de temps \(\tau=RC\) représente l'\textbf{échelle temporelle caractéristique}
du phénomène de charge.
\end{imagem}

\begin{erreurs}
\begin{itemize}
    \item confondre \(u_R\) et \(u_C\) ;
    \item oublier la loi des mailles ;
    \item ne pas relier l'exponentielle au comportement graphique.
\end{itemize}
\end{erreurs}

\newpage

\section{Terminale générale -- Ondes et signaux}

\subsection{Lentilles et optique}

\begin{sens}
L'optique géométrique permet de modéliser la formation des images
à partir de rayons lumineux et de lois simples.
\end{sens}

\begin{imagem}
Une lentille convergente peut être vue comme un \textbf{organisateur de rayons}
qui fait converger ou diverger la lumière selon la situation.
\end{imagem}

\begin{erreurs}
\begin{itemize}
    \item confusion entre objet réel et image ;
    \item schémas de rayons faux ;
    \item mauvais usage des signes ;
    \item confusion entre taille géométrique et netteté.
\end{itemize}
\end{erreurs}

\subsection{Effet Doppler}

\begin{sens}
L'effet Doppler montre que la fréquence perçue dépend du mouvement relatif de la source et du récepteur.
\end{sens}

\begin{imagem}
Quand la source se rapproche, les ondes arrivent \textbf{plus serrées} ;
quand elle s'éloigne, elles arrivent \textbf{plus espacées}.
\end{imagem}

\begin{fausseintuition}
Les élèves peuvent croire que la source émet réellement une fréquence différente.
Il faut préciser que ce qui change est la fréquence reçue ou observée.
\end{fausseintuition}

\subsection{Interférences, diffraction, aspects ondulatoires}

\begin{sens}
Ces phénomènes montrent que la lumière et d'autres signaux ondulatoires
ne se comportent pas seulement comme des trajectoires rectilignes simples.
\end{sens}

\begin{imagem}
Une onde peut se répartir, se recombiner, produire des renforcements ou des annulations :
elle se comporte comme une \textbf{structure étendue} et non comme un simple projectile.
\end{imagem}

\begin{idee}
Ces chapitres sont particulièrement utiles pour faire comprendre que les modèles physiques changent selon le phénomène étudié.
\end{idee}

\newpage

\section{Mesures, incertitudes et traitement des données}

\subsection{Mesure}

\begin{sens}
Une mesure n'est jamais une vérité absolue :
c'est une estimation associée à une précision limitée.
\end{sens}

\begin{imagem}
Mesurer, ce n'est pas \og trouver la vraie valeur \fg,
c'est \textbf{encadrer raisonnablement une valeur plausible}.
\end{imagem}

\begin{fausseintuition}
Les élèves prennent souvent la valeur affichée par un appareil comme une vérité parfaite.
\end{fausseintuition}

\subsection{Incertitude}

\begin{sens}
L'incertitude traduit le manque de précision inévitable associé à une mesure ou à un résultat.
\end{sens}

\begin{imagem}
L'incertitude est la \textbf{marge raisonnable de doute} autour d'une valeur mesurée.
\end{imagem}

\begin{fausseintuition}
Les élèves pensent parfois qu'une incertitude signifie que la mesure est mauvaise.
Au contraire, donner une incertitude est une marque de rigueur scientifique.
\end{fausseintuition}

\begin{erreurs}
\begin{itemize}
    \item trop de chiffres significatifs ;
    \item arrondis incohérents ;
    \item comparaison de valeurs sans tenir compte des incertitudes ;
    \item oubli de l'unité.
\end{itemize}
\end{erreurs}

\newpage

\section{Grandes difficultés conceptuelles transversales}

\begin{longtable}{p{0.28\textwidth} p{0.67\textwidth}}
\toprule
\textbf{Confusion fréquente} & \textbf{Comment la lever} \\
\midrule
Force / mouvement & Répéter qu'une force modifie le mouvement, elle n'est pas synonyme de mouvement. \\
Vitesse / accélération & Multiplier les exemples où la vitesse est grande mais l'accélération nulle, et inversement. \\
Masse / poids & Revenir au fait que la masse mesure l'inertie et que le poids est une force. \\
Énergie / puissance & Dire systématiquement : énergie = quantité ; puissance = rythme. \\
Température / chaleur & Opposer état thermique et transfert d'énergie. \\
Charge / courant & Dire : charge = quantité électrique ; courant = débit de charge. \\
Tension / courant & La tension compare deux points ; le courant traverse une section. \\
Onde / transport de matière & Faire observer qu'une perturbation se propage sans transporter globalement le milieu. \\
Quantité de matière / masse & Revenir à l'idée de comptage d'entités. \\
Modèle / réalité & Poser à chaque fois la question : qu'a-t-on négligé ? \\
\bottomrule
\end{longtable}

\section{Erreurs classiques de méthode}

\subsection{En mécanique}
\begin{itemize}
    \item oublier de définir le système étudié ;
    \item ne pas préciser le référentiel ;
    \item faire un bilan des forces incomplet ;
    \item écrire une relation sans projeter ;
    \item ne pas vérifier le signe ou l'ordre de grandeur du résultat.
\end{itemize}

\subsection{En chimie}
\begin{itemize}
    \item ne pas équilibrer l'équation-bilan ;
    \item mal identifier les espèces en jeu ;
    \item oublier les unités ;
    \item utiliser la stœchiométrie sans la justifier ;
    \item ne pas commenter le résultat obtenu.
\end{itemize}

\subsection{En électricité}
\begin{itemize}
    \item confondre tension, courant, énergie et puissance ;
    \item mal utiliser les conventions de signe ;
    \item ne pas exploiter le schéma global du montage avant le calcul.
\end{itemize}

\subsection{En ondes et optique}
\begin{itemize}
    \item schémas imprécis ;
    \item axes non légendés ;
    \item confusion entre ce qui oscille localement et ce qui se propage.
\end{itemize}

\section{Répertoire des démonstrations importantes à maîtriser}

\begin{itemize}
    \item relation de dilution ;
    \item structure du tableau d'avancement ;
    \item principe d'inertie ;
    \item forme projetée de la deuxième loi de Newton ;
    \item équations horaires de la chute libre ;
    \item expression du travail d'une force constante ;
    \item théorème de l'énergie cinétique ;
    \item conservation de l'énergie mécanique ;
    \item relation \(f=1/T\) ;
    \item loi exponentielle \(N(t)=N_0\e^{-\lambda t}\) ;
    \item relation à l'équivalence dans un titrage ;
    \item équation différentielle d'un circuit RC ;
    \item vitesse orbitale \(v=\sqrt{GM/r}\).
\end{itemize}

\section{Conseils pédagogiques pour l'enseignement}

\begin{enumerate}
    \item Commencer par une situation concrète ou une expérience.
    \item Faire émerger les grandeurs pertinentes.
    \item Expliquer le choix du modèle utilisé.
    \item Distinguer clairement observation, hypothèse, loi, calcul et interprétation.
    \item Faire verbaliser le sens des grandeurs avant l'application des formules.
    \item Revenir régulièrement sur les fausses intuitions.
    \item Demander systématiquement une vérification finale :
    \begin{itemize}
        \item unité ;
        \item signe ;
        \item cohérence physique ;
        \item ordre de grandeur.
    \end{itemize}
\end{enumerate}

\begin{idee}
Quelques questions simples à poser très souvent en classe :
\begin{itemize}
    \item Quel est le système étudié ?
    \item Dans quel référentiel travaille-t-on ?
    \item Quelles interactions garde-t-on ? Quelles interactions néglige-t-on ?
    \item Que mesure réellement cette grandeur ?
    \item Ce résultat a-t-il du sens physiquement ?
\end{itemize}
\end{idee}

\section{Conclusion}

\begin{bilan}
Le coeur de l'enseignement de la physique-chimie au lycée n'est pas l'application de recettes,
mais la construction d'une manière scientifique de penser :
\begin{itemize}
    \item observer ;
    \item modéliser ;
    \item mesurer ;
    \item calculer ;
    \item interpréter ;
    \item critiquer.
\end{itemize}

Une bonne progression pédagogique consiste à faire apparaître à chaque chapitre :
\begin{enumerate}
    \item ce qu'on cherche à décrire ;
    \item avec quelle grandeur ;
    \item avec quel modèle ;
    \item pourquoi ce modèle est raisonnable ici ;
    \item quelles erreurs de sens il faut éviter.
\end{enumerate}
\end{bilan}

\end{document}